\begin{textsample}{POS Dim 5 – global\_north\_2025\_02 – Score 12.00 – global\_north\_2025\_02/121237575.txt}  \label{ex:f5_pos_336}
The wheels of justice may turn slowly, but they ' re finally catching up with officials at the \textbf{International} Criminal \textbf{Court} after President Donald Trump signed an executive order \textbf{sanctioning} those who target U.S. citizens or allies.
This comes on the heels of the House-passed Illegitimate \textbf{Court} Counteraction \textbf{Act}.
And while the Senate failed to pass its version of the bill, Senate Minority Leader Chuck Schumer signaled the Democrats are aligned with the president ' s approach.
These \textbf{sanctions} are a good start, but to truly curb the ICC ' s harmful actions, the U.S. will need to convince its allies to follow suit.
The impetus for the U.S. action is the ICC ' s recent charges against Israel, an examination of which shows the body is guilty of abusing its \textbf{legal} power -- selectively choosing data and skipping due process.
The attack on Israel escalated on Nov. 21, when the ICC issued arrest \textbf{warrants} for Israel ' s prime minister and former minister of defense, charging them with committing"the war \textbf{crime} of starvation as a \textbf{method} of @ @ @ @ @ @ @ @ @ @ to May 20.
The ICC ' s use of Oct. 8 as a start date would have been comical if it weren't so tragic.
It picked one day after Israel suffered the worst \textbf{genocidal} assault on Jews since the Holocaust -- before Israel ' s military even entered Gaza.
As for the charge itself, my analysis of data easily available to the ICC -- but which it chose to ignore -- demonstrates that far from"impeding humanitarian aid,"Israel has gone to great pains to facilitate the delivery of aid and food since the war ' s start.
Using comprehensive data tracked by COGAT ( Israel ' s Coordinator of Government Activities in the Territories ), I found that between Nov. 1, 2023, and May 20, an average of 2,152 tons of food was brought into Gaza per day.
Full distribution of this amount to the 2.1 million people in Gaza would have supplied each Gazan with 1.02 kilograms of food per day ( 1 kilogram of food is typically equivalent to more than 3,000 calories ). @ @ @ @ @ @ @ @ @ @ 94,000 metric tons of food are"enough to feed one million people for four months,"which translates to 0.78 kilograms of food per person per day.
Thus, using the U.N. ' s own standards, the food supplied surpassed 130 \% of the population ' s daily dietary needs.
Although COGAT is a party to the conflict, it engages in coordinated initiatives with U.N. agencies, nongovernmental organizations and the U.S.-supported Joint Coordination Board, and its data is reviewed by representatives from the U.S., Egypt and the U.N.
I interviewed Maj.
Gen.
Ghassan Alian, the head of COGAT ; he told me this coordination network was set up to refute unsubstantiated \textbf{allegations} of land crossing closures and insufficient food supply.
The ICC and U.N., however, disfavored the COGAT data, relying instead on incomplete data and biased analyses.
For example, the U.N.
Office for the Coordination of Humanitarian Affairs disclosed in an April report that the underlying data it used, collected by the U.N.
Relief and Works Agency for @ @ @ @ @ @ @ @ @ @ delivered via trucks through land crossings at Kerem Shalom and Rafah but not others, such as the Erez crossing in northern Gaza.
And it included only aid observed and registered by UNRWA representatives -- when they were on duty.
By contrast, COGAT reports on"all aid that enters the Gaza Strip, from all sources and through all crossings,"according to the Israeli think tank Institute for National Security Studies.
Behavioral economics and human perception -- two of my areas of academic research -- help explain why the ICC ' s false \textbf{accusations} are so harmful.
For example, simply repeating a message increases the likelihood of a recipient both recalling that message and perceiving it as true -- even if that content is implausible.
Once false narratives against Israel were propagated, they became dutifully repeated.
Research has also shown that messages framed in a more negative, extreme or visual manner are more attention-grabbing and are more likely to be shared with others.
America has a track record of not standing idly by while @ @ @ @ @ @ @ @ @ @ understanding on both sides of the political aisle that, left unopposed, the ICC could similarly try to curtail the West ' s ability to defend itself.
This is why the U.S. has long opposed the ICC and its claims of \textbf{legal} \textbf{jurisdiction}.
In fact, the ICC has tried to target the U.S. before -- authorizing an investigation into alleged war \textbf{crimes} committed in Afghanistan by U.S. troops -- but U.S. pressure curtailed it.
The reality, however, is that U.S. \textbf{sanctions} on the ICC are unlikely to be enough.
The longer-term goal needs to be to shut down the \textbf{court} altogether or fundamentally reform it.
This will likely require intense negotiations with countries who fund it -- which include U.S. allies such as France, Canada and the United Kingdom.
Today ' s U.S. isn't afraid of tough conversations with allies, including pushing them to spend more on their militaries as part of their NATO commitment.
Now that ' s a far better use of their ICC funds.

% matched lemmas: accusation, act, allegation, court, crime, genocidal, international, jurisdiction, legal, method, sanction, warrant
\end{textsample}
